\section{Class 3: Internal Analysis}
Resource Based View (RBV): Competitive advantage arises from unique, valuable,
and hard-to-imitate resources. Focus on leveraging internal strengths to
exploit external opportunities and mitigate threats.

Has to be non-fungible, valuable, rare, and hard to imitate. Examples: Brand
equity (Starbucks),Hub dominance (Airlines), and R\&D capabilities (IBM). These
resources enable firms to create superior value and sustain competitive
advantage.

\subsection{VRIO Framework}
\textbf{Valuable:} Does the resource enable the firm to exploit opportunities or
neutralize threats? If not, it may be a competitive disadvantage.

\textbf{Rare:} Is the resource controlled by a few firms? If not, it may lead to
competitive parity.

\textbf{Inimitable:} Is the resource difficult to imitate due to unique historical
conditions, causal ambiguity, or social complexity? If not, it may lead to
temporary competitive advantage.

\textbf{Organized:} Is the firm organized to capture value from the resource? If not,
it may lead to unrealized competitive advantage.

\begin{tikzpicture}[
        font=\sffamily\tiny,
        >=Stealth,
        node distance=0.5cm and 0.8cm,
        diamond_node/.style={draw, diamond, aspect=1, minimum size=1.8cm, thick, align=center, text=white},
        result_node/.style={draw=blue!30, fill=blue!10, rounded corners=3pt, minimum width=2cm, minimum height=1cm, align=center, thick},
        label_node/.style={font=\sffamily\bfseries\tiny, align=center},
        scale=0.7, transform shape
    ]

    % Top Labels
    \node[label_node] (res_label) at (1, 1.5) {IS THE RESOURCE OR CAPABILITY...?};
    \node[label_node] (comp_label) at (5, 1.5) {IS THE COMPANY WELL...?};

    % Diamonds
    \node[diamond_node, fill=orange] (v) {\textbf{V} \\ VALUABLE}; \node[diamond_node, fill=purple!80!black, right=of v] (r) {\textbf{R} \\ RARE}; \node[diamond_node, fill=green!50!black, right=of r] (i) {\textbf{I} \\ INIMITABLE}; \node[diamond_node, fill=cyan!70!black, right=of i] (o) {\textbf{O} \\ ORGANIZED};

    % Results Bottom
    \node[result_node, below=0.8cm of v] (res_v) {COMPETITIVE \\ DISADVANTAGE};
    \node[result_node, below=0.8cm of r] (res_r) {COMPETITIVE \\ PARITY};
    \node[result_node, below=0.8cm of i] (res_i) {TEMPORARY \\ COMPETITIVE \\ ADVANTAGE};
    \node[result_node, below=0.8cm of o] (res_o) {UNUSED \\ COMPETITIVE \\ ADVANTAGE};

    % Final Result Right
    \node[result_node, right=0.8cm of o, fill=blue!20] (res_final) {SUSTAINED \\ COMPETITIVE \\ ADVANTAGE};

    % Paths with "YES" and "NO"
    \path[->, thick] (v) edge node[above] {YES} (r);
    \path[->, thick] (r) edge node[above] {YES} (i);
    \path[->, thick] (i) edge node[above] {YES} (o);
    \path[->, thick] (o) edge node[above] {YES} (res_final);

    \path[->, thick] (v) edge node[left] {NO} (res_v);
    \path[->, thick] (r) edge node[left] {NO} (res_r);
    \path[->, thick] (i) edge node[left] {NO} (res_i);
    \path[->, thick] (o) edge node[left] {NO} (res_o);

\end{tikzpicture}

\subsubsection{Building or Buying Resources}
\textbf{Building:} Develop resources internally through R\&D, training, or culture.
Takes time and may require significant investment but can create unique
capabilities.

\textbf{Buying:} Acquire resources through mergers, acquisitions, or partnerships.
Faster but may be costly and risk integration issues. Can be effective for
acquiring complementary resources or entering new markets.

\begin{callout}
    Successful M\&A integrates activities that are related yet complementary,
    creating synergies that enhance competitive advantage. For example, Disney
    acquiring Pixar combined Disney's storytelling and distribution with Pixar's
    cutting-edge animation technology, resulting in a stronger competitive position
    in the animation industry.
\end{callout}

\section{Firm Growth Theory}
Only expand into a market or start a new activity if you can leverage your core
resources and capabilities to create a competitive advantage. This is the basis
for the \textbf{Related Diversification} strategy, which focuses on entering markets or
activities that are related to the firm's existing operations, allowing it to
leverage its core competencies and achieve synergies.

In contrast, \textbf{Unrelated Diversification} involves entering markets or activities that are not
related to the firm's existing operations, which can be riskier and may not
leverage the firm's core competencies effectively. Turning into a conglomerate
can lead to inefficiencies and lack of focus, as the firm may struggle to
manage diverse businesses effectively.

